\documentclass[aspectratio=169,10pt]{beamer}
\usetheme[numbering=fraction,progressbar=frametitle]{metropolis}
\usepackage{fontspec}
\setsansfont{Latin Modern Sans}
\usepackage{amsmath,amssymb,mathtools,graphicx,booktabs}
\usepackage{polyglossia}
\setdefaultlanguage{french}
\setbeamertemplate{navigation symbols}{}
\setbeamerfont{frametitle}{size=\large}
\setbeamerfont{normal text}{size=\normalsize}
\setlength{\parskip}{.45em}
\hypersetup{pdfauthor={},pdfsubject={Lecture slides},pdfkeywords={},pdftitle={Géométrie dans l’espace}}
\providecommand{\N}{\mathbb{N}}
\providecommand{\Z}{\mathbb{Z}}
\providecommand{\D}{\mathbb{D}}
\providecommand{\Q}{\mathbb{Q}}
\providecommand{\R}{\mathbb{R}}
\providecommand{\C}{\mathbb{C}}
\providecommand{\vect}[1]{\overrightarrow{#1}}
\providecommand{\vi}{\vec{\imath}}
\providecommand{\vj}{\vec{\jmath}}
\providecommand{\vk}{\vec{k}}
\providecommand{\norm}[1]{\left\lVert #1\right\rVert}
\providecommand{\vc}[2]{\begin{pmatrix}#1\\#2\end{pmatrix}}
\providecommand{\deter}[2]{\det\left(#1,#2\right)}
\providecommand{\abs}[1]{\left\lvert #1\right\rvert}
\providecommand{\dd}{\,\mathrm{d}}
\providecommand{\e}{\mathrm{e}}
\providecommand{\Card}{\operatorname{Card}}
\providecommand{\E}{\operatorname{E}}
\providecommand{\V}{\operatorname{V}}
\providecommand{\degre}{^{\circ}}
\providecommand{\un}[1]{\,\mathrm{#1}}
\providecommand{\sci}[2]{#1\times10^{#2}}
\providecommand{\pgcd}{\operatorname{PGCD}}
\providecommand{\ppcm}{\operatorname{PPCM}}
\title{Géométrie dans l’espace}
\subtitle{Cours et exemples}
\date{}
\author{}
\begin{document}
\begin{frame}[plain]
\vfill
{\usebeamerfont{title}\inserttitle\par}
\medskip
{\usebeamerfont{subtitle}\insertsubtitle\par}
\medskip
{\usebeamercolor[fg]{progress bar}\rule{\textwidth}{.5pt}}
\vfill
\end{frame}
\begin{frame}{Objectifs}
Utiliser un repère direct, calculer produits scalaire et vectoriel et établir les équations des droites et des plans.
\end{frame}
\begin{frame}{Repérage et orientation}
Un repère de l’espace comporte une origine et trois vecteurs de base non coplanaires. Un repère orthonormé utilise trois vecteurs unitaires orthogonaux. Un repère orthonormé est direct lorsque le produit vectoriel du premier et du deuxième vecteur de base donne le troisième. Dans une base quelconque, être direct exprime seulement une orientation, pas cette égalité.
\end{frame}
\begin{frame}{Repérage et orientation}
On repère un cube par \(O(0;0;0)\), \(A(1;0;0)\), \(B(0;1;0)\), \(C(0;0;1)\). Le sommet opposé à \(O\) a coordonnées \((1;1;1)\) et la grande diagonale vaut \(\sqrt3\). Une perspective cavalière aide à lire la configuration, mais ses longueurs dessinées ne donnent pas les longueurs de l’espace.
\end{frame}
\begin{frame}{Repérage et orientation}
Dans un repère orthonormé, pour \(\overrightarrow{u}=(u_1;u_2;u_3)\) et \(\overrightarrow{v}=(v_1;v_2;v_3)\), \(\overrightarrow{u}\cdot\overrightarrow{v}=u_1v_1+u_2v_2+u_3v_3\). La norme vaut \(\sqrt{u_1^2+u_2^2+u_3^2}\), et, pour deux vecteurs non nuls, le produit géométrique vaut \(\left\| \overrightarrow{u} \right\|\left\| \overrightarrow{v} \right\|\cos\theta\); si l’un est nul, le produit vaut zéro.
\end{frame}
\begin{frame}{Produit vectoriel}
\[
\overrightarrow{u}\wedge\overrightarrow{v}=(u_2v_3-u_3v_2\,;\,
 u_3v_1-u_1v_3\,;\,u_1v_2-u_2v_1).
\]
Ce vecteur est orthogonal aux deux vecteurs, et son sens dépend de l’ordre. Si les deux vecteurs sont non colinéaires, le triplet \((\overrightarrow{u},\overrightarrow{v},\overrightarrow{u}\wedge\overrightarrow{v})\) est direct.
\end{frame}
\begin{frame}{Produit vectoriel}
\(\overrightarrow{v}\wedge\overrightarrow{u}=-\overrightarrow{u}\wedge\overrightarrow{v}\); pour deux vecteurs non nuls, sa norme vaut \(\left\| \overrightarrow{u} \right\|\left\| \overrightarrow{v} \right\|\sin\theta\) avec \(\theta\in[0;\pi]\); si l’un est nul, le produit vectoriel est nul. Elle est l’aire du parallélogramme, donc la moitié donne l’aire du triangle. Le produit est nul exactement lorsque les vecteurs sont colinéaires.
\end{frame}
\begin{frame}{Produit vectoriel}
Le produit scalaire avec chaque facteur s’annule en développant les coordonnées. Développer son carré donne \(\left\| \overrightarrow{u}\wedge\overrightarrow{v} \right\|^2=\left\| \overrightarrow{u} \right\|^2\left\| \overrightarrow{v} \right\|^2-(\overrightarrow{u}\cdot\overrightarrow{v})^2\), soit la formule avec \(\sin^2\theta\). Le sinus est positif ou nul dans \([0;\pi]\). L’orientation relève de la convention du repère direct.
\end{frame}
\begin{frame}{Produit vectoriel}
Pour \(\overrightarrow{u}=(1;2;0)\) et \(\overrightarrow{v}=(0;1;3)\), \(\overrightarrow{u}\wedge\overrightarrow{v}=(6;-3;1)\). Produits scalaires \(6-6=0\) et \(-3+3=0\). Norme \(\sqrt{46}\), aire du triangle \(\sqrt{46}/2\).
\end{frame}
\begin{frame}{Droites et plans}
La droite passant par \(A(a;b;c)\) de vecteur directeur non nul \((u;v;w)\) a pour équations paramétriques \(x=a+ut\), \(y=b+vt\), \(z=c+wt\), \(t\in\mathbb{R}\). Dans l’espace, une seule équation cartésienne linéaire ne décrit pas une droite; on utilise un système de deux équations de plans non parallèles.
\end{frame}
\begin{frame}{Droites et plans}
Un plan passant par \(A\) et dirigé par deux vecteurs non colinéaires \(\overrightarrow{u},\overrightarrow{v}\) est décrit par \(\overrightarrow{AM}=s\overrightarrow{u}+t\overrightarrow{v}\), \(s,t\in\mathbb{R}\). Pour un vecteur normal non nul \(\overrightarrow{n}=(\alpha;\beta;\gamma)\), il a une équation \(\alpha x+\beta y+\gamma z+\delta=0\).
\end{frame}
\begin{frame}{Droites et plans}
L’équation normale résulte de \(\overrightarrow{n}\cdot\overrightarrow{AM}=0\), avec \(\delta=-\alpha a-\beta b-\gamma c\). Pour trois points non alignés, leur produit vectoriel fournit un vecteur normal.
\end{frame}
\begin{frame}{Droites et plans}
Substituer la représentation paramétrique de la droite dans l’équation du plan. Une équation linéaire non constante donne un point unique; une égalité impossible donne aucun point (parallélisme strict); une identité donne la droite incluse dans le plan.
\end{frame}
\begin{frame}{Droites et plans}
Pour \(P:\alpha x+\beta y+\gamma z+\delta=0\), la distance de \(A(a;b;c)\) à \(P\) vaut \(|\alpha a+\beta b+\gamma c+\delta|/\sqrt{\alpha^2+\beta^2+\gamma^2}\).
\end{frame}
\begin{frame}{Droites et plans}
La projection orthogonale \(H\) de \(A\) vérifie \(\overrightarrow{HA}=\lambda\overrightarrow{n}\). Remplacer dans l’équation du plan donne la longueur \(AH\) égale à la formule. Le vecteur normal non nul permet la division.
\end{frame}
\begin{frame}{Droites et plans}
Trouver une normale non nulle, par exemple un produit vectoriel de deux directions non colinéaires. Écrire \(\overrightarrow{n}\cdot\overrightarrow{AM}=0\), puis développer. Vérifier les points de départ dans l’équation finale. Pour une droite, utiliser au contraire un vecteur directeur dans la représentation paramétrique; le normal d’un plan n’est pas une direction contenue dans ce plan.
\end{frame}
\begin{frame}{Droites et plans}
Pour \(P:x+2y-z=3\) et \(A(1;0;0)\), la normale est \(\overrightarrow{n}=(1;2;-1)\), de norme au carré \(6\). La valeur du membre \(x+2y-z-3\) en \(A\) est \(-2\). Donc
\[
H=A-\dfrac{-2}{6}\overrightarrow{n}=(4/3;2/3;-1/3).
\]
Sa substitution donne \(3\); le vecteur \(\overrightarrow{AH}\) est colinéaire à \(\overrightarrow{n}\). Le point est bien le projeté orthogonal, et \(AH=\sqrt6/3\), ce que confirme la formule de distance \(2/\sqrt6\).
\end{frame}
\begin{frame}{Droites et plans}
Des normales proportionnelles indiquent des plans parallèles, éventuellement confondus. Comparer aussi la constante : multiplier toute l’équation par un même réel non nul ne change pas le plan. Si les normales ne sont pas colinéaires, résoudre les deux équations pour obtenir leur droite d’intersection.
\end{frame}
\begin{frame}{Droites et plans}
Une substitution donnant \(0=0\) signifie que tous les paramètres conviennent, pas que le paramètre vaut zéro. Une substitution impossible interdit tous les points de la droite. Ces trois cas sont les analogues géométriques des issues d’un système linéaire.
\end{frame}
\end{document}
